%! Functions — Unit 1, Lesson 1.3 — Warm-Up
\documentclass[10pt]{article}
\usepackage{algebra2-article}
\usepackage{algebra2-boxes}
\newcommand{\TallMath}[1]{$\displaystyle #1\rule[-1.4em]{0pt}{3.2em}$}

% Pre-drawn coordinate grid; #1 receives the plotted curve and any overlay.
\newcommand{\gridplot}[1]{%
  \begin{tikzpicture}[x=0.36cm, y=0.24cm]
    \draw[step=1, linegray!45, very thin] (-5,-7) grid (5,7);
    \draw[->, charcoal, thick] (-5.4,0) -- (5.4,0) node[right, font=\scriptsize] {$x$};
    \draw[->, charcoal, thick] (0,-7.4) -- (0,7.4) node[above, font=\scriptsize] {$y$};
    \foreach \t in {-4,-2,2,4} {\node[below right=-2pt, font=\tiny] at (\t,0) {$\t$};}
    \foreach \t in {-6,-4,-2,2,4,6} {\node[above left=-2pt, font=\tiny] at (0,\t) {$\t$};}
    #1
  \end{tikzpicture}}

\begin{document}

\pageheader{Unit 1, Lesson 1.3}{Warm-Up}
\namedateperiod

\vspace{0.06in}

\begin{notesbox}{Getting Ready: Three Ways to Ask One Question}
\small Nothing here is new. Work quickly --- the point is what the three items have in common.

\vspace{3pt}
\textbf{1. Solve (Lesson 1.2).} Find \emph{every} value that makes each expression equal zero.
\begin{center}
\renewcommand{\arraystretch}{1.6}
\begin{tabular}{@{}l l@{}}
(a) \; $2x-6=0$, \quad $x=$ \blank{3.0cm} &
(b) \; $x^2-2x-3=0$, \quad $x=$ \blank{3.6cm} \\
\end{tabular}
\end{center}

\vspace{2pt}
\textbf{2. Evaluate (Lesson 1.1).} Each rule below is a recipe. Run it at the two inputs given.
\begin{center}
\renewcommand{\arraystretch}{1.45}
\begin{tabularx}{\linewidth}{@{}X c c@{}}
\hline
\textbf{Rule} & \textbf{Value at $x=0$} & \textbf{Value at $x=3$} \\
\hline
$2x-6$ & \blank{2.4cm} & \blank{2.4cm} \\
$x^2-2x-3$ & \blank{2.4cm} & \blank{2.4cm} \\
$2^x$ & \blank{2.4cm} & \blank{2.4cm} \\
\hline
\end{tabularx}
\end{center}
Which rules gave you \textbf{zero} at $x=3$? \blank{4.6cm} \\
Now look back at item 1. Is that a coincidence? \blank{2.0cm} \\
Is there \emph{any} input that makes $2^x$ equal zero? \blank{1.6cm} \; (Try a few. Try a negative one.)

\vspace{3pt}
\textbf{3. Read the picture.} Here is a graph. Answer only from the picture.
\begin{center}
\begin{tabular}{@{}c l@{}}
\gridplot{\draw[forest, very thick, domain=-0.6:5, smooth, variable=\x] plot ({\x},{2*\x-6});}
&
\begin{minipage}[c]{0.52\linewidth}
It crosses the $x$-axis at $x=$ \blank{1.6cm}
\\[6pt]
It crosses the $y$-axis at $y=$ \blank{1.6cm}
\\[6pt]
Left to right, it goes \blank{2.4cm} (up / down).
\\[6pt]
Every time $x$ increases by $1$, $y$ changes by \blank{1.4cm}
\end{minipage}
\\
\end{tabular}
\end{center}

\vspace{1pt}
\textbf{Put it together.} Item 1(a), the first row of item 2, and the graph above are all the
\emph{same rule}. How can you tell? \blank{8.0cm}

\end{notesbox}

\end{document}
